\chapter{Immunotherapy}

\section*{Overview}
Immunotherapy stimulates or trains the body's own immune system to recognize and destroy cancer cells. This includes immune checkpoint inhibitors, monoclonal antibodies, and T-cell therapies.

\section*{Alternative Names}
Biological therapy, Immuno-oncology

\section*{Principle}
Treatments are used to stimulate, suppress, or modulate the immune system: allergen immunotherapy induces tolerance, checkpoint inhibitors release antitumor T-cell responses, and biologic agents target inflammatory pathways. The immune response against the disease is thereby altered.
\section*{History}
William Coley used bacterial toxins ('Coley's toxins') in the 1890s to stimulate the immune system. James Allison and Tasuku Honjo won the 2018 Nobel Prize for discovering immune checkpoint inhibition.

\section*{Why This Test or Procedure Is Ordered}
Ordered to treat various advanced cancers, including melanoma, lung cancer, kidney cancer, and lymphoma, often when standard therapies have failed.

\section*{Is This a Primary or Secondary Test or Procedure}
Primary systemic treatment for advanced melanoma and specific non-small cell lung cancers.

\section*{How This Test or Procedure Is Performed}
Usually administered intravenously in an outpatient clinic. Infusions occur every few weeks in cycles, similar to chemotherapy.

\section*{What Preparation Is Needed}
Complete blood count, metabolic panel, and thyroid function tests are performed prior to each infusion to screen for immune-related adverse events.

\section*{Contraindications and Precautions}
Active severe autoimmune diseases (e.g., lupus, Crohn's) are relative contraindications due to the risk of exacerbating the disease.

\section*{Risks}
Immune-related adverse events (irAEs) including colitis, pneumonitis, hepatitis, thyroiditis, hypophysitis, or severe skin rashes.

\section*{What Happens After the Test or Procedure}
Immediately report any new symptoms of shortness of breath, severe diarrhea, abdominal pain, or extreme fatigue to your oncologist.

\section*{Understanding Your Results}
Assessed using imaging scans. Response may take longer than chemotherapy, and tumors may temporarily appear to grow before shrinking (pseudoprogression).

\section*{Normal Results}
Tumor response or disease stability; absence of severe immune-mediated organ toxicity.

\section*{Follow-up Tests and Procedures}
Frequent monitoring of thyroid function, liver enzymes, and clinical symptoms.

\section*{References}
\begin{enumerate}
  \item MedlinePlus. Immunotherapy. U.S. National Library of Medicine; 2025.
  \item Current Medical Diagnosis and Treatment. McGraw-Hill; 2025.
\end{enumerate}

\clearpage
